\documentclass[11pt]{article}
\usepackage{mystyle}

\title{Rosen --- \S 1.5: Nested Quantifiers\\ \large Selected Exercises}
\author{Alston}
\date{Semester 2, 2026}

\begin{document}
\maketitle


% ======================================================================
\begin{exercise}{1}
    Translate these statements into English, where the domain for
    each variable consists of all real numbers.
    \begin{enumerate}[label=(\alph*)]
        \item $\forall x\,\exists y\,(x<y)$
        \item $\forall x\,\forall y\,\big(((x\geq 0)\wedge(y\geq 0))\to(xy\geq 0)\big)$
        \item $\forall x\,\forall y\,\exists z\,(xy=z)$
    \end{enumerate}
\end{exercise}

\begin{solution}
\textbf{(a)} $\forall x\,\exists y\,(x<y)$

For all $x$, there exists $y$ such that $x<y$. (For every real
number, there is a larger real number.)

\bigskip
\noindent\textbf{(b)} $\forall x\,\forall y\,\big((x\geq0\wedge y\geq0)\Rightarrow xy\geq0\big)$

For all real numbers $x$ and $y$, if $x$ and $y$ are both
nonnegative, then their product $xy$ is nonnegative. (The product of
two nonnegative real numbers is nonnegative.)

\bigskip
\noindent\textbf{(c)} $\forall x\,\forall y\,\exists z\,(xy=z)$

For all $x$ and $y$, there exists $z$ such that $xy=z$. (Every pair
of real numbers has a product.)
\end{solution}

% ======================================================================
\begin{exercise}{7}
    Let $T(x,y)$ mean that student $x$ likes cuisine $y$, where the
    domain for $x$ consists of all students at your school and the
    domain for $y$ consists of all cuisines. Express each of these
    statements by a simple English sentence.
    \begin{enumerate}[label=(\alph*)]
        \item $\neg T(\text{Abdallah Hussein}, \text{Japanese})$
        \item $\exists x\,T(x,\text{Korean})\wedge \forall x\,T(x,\text{Mexican})$
        \item $\exists y\,\big(T(\text{Monique Arsenault}, y)\vee T(\text{Jay Johnson}, y)\big)$
        \item $\forall x\,\forall z\,\exists y\,\big((x\neq z)\to \neg(T(x,y)\wedge T(z,y))\big)$
        \item $\exists x\,\exists z\,\forall y\,\big(T(x,y)\leftrightarrow T(z,y)\big)$
        \item $\forall x\,\forall z\,\exists y\,\big(T(x,y)\leftrightarrow T(z,y)\big)$
    \end{enumerate}
\end{exercise}

\begin{solution}
Let $T(x,y)$ mean ``student $x$ likes cuisine $y$.''

\medskip
\noindent\textbf{(a)} $\neg T(\text{Abdallah Hussein}, \text{Japanese})$

Abdallah Hussein does not like Japanese cuisine.

\medskip
\noindent\textbf{(b)} $\exists x\,T(x,\text{Korean})\wedge \forall x\,T(x,\text{Mexican})$

There is a student who likes Korean cuisine, and every student likes
Mexican cuisine.

\medskip
\noindent\textbf{(c)} $\exists y\,\big(T(\text{Monique Arsenault}, y)\vee T(\text{Jay Johnson}, y)\big)$

There is a cuisine that Monique Arsenault or Jay Johnson likes.

\medskip
\noindent\textbf{(d)} $\forall x\,\forall z\,\exists y\,\big((x\neq z)\to \neg(T(x,y)\wedge T(z,y))\big)$

For any two different students, there is some cuisine that they
don't both like.

\medskip
\noindent\textbf{(e)} $\exists x\,\exists z\,\forall y\,\big(T(x,y)\leftrightarrow T(z,y)\big)$

There exist two students $x$ and $z$ such that, for every cuisine
$y$, $x$ likes $y$ if and only if $z$ likes $y$. (There are two
students with exactly the same cuisine preferences.)

\medskip
\noindent\textbf{(f)} $\forall x\,\forall z\,\exists y\,\big(T(x,y)\leftrightarrow T(z,y)\big)$

For all students $x$ and $z$, there is some cuisine $y$ such that
$x$ likes $y$ if and only if $z$ likes $y$. (This is trivially true
for every pair of students, e.g.\ any cuisine both like or both
dislike.)
\end{solution}

% ======================================================================
\begin{exercise}{9}
    Let $L(x,y)$ be the statement ``$x$ loves $y$,'' where the domain
    for both $x$ and $y$ consists of all people in the world. Use
    quantifiers to express each of these statements.
    \begin{enumerate}[label=(\alph*)]
        \item Everybody loves Jerry.
        \item Everybody loves somebody.
        \item There is somebody whom everybody loves.
        \item Nobody loves everybody.
        \item There is somebody whom Lydia does not love.
        \item There is somebody whom no one loves.
        \item There is exactly one person whom everybody loves.
        \item There are exactly two people whom Lynn loves.
        \item Everyone loves himself or herself.
        \item There is someone who loves no one besides himself or herself.
    \end{enumerate}
\end{exercise}

\begin{solution}
Let $L(x,y)$ mean ``$x$ loves $y$,'' with domain all people.

\medskip
\noindent\textbf{(a)}
\[
\forall x\, L(x,\text{Jerry})
\]

\medskip
\noindent\textbf{(b)}
\[
\forall x\,\exists y\, L(x,y)
\]

\medskip
\noindent\textbf{(c)}
\[
\exists y\,\forall x\, L(x,y)
\]

\medskip
\noindent\textbf{(d)}
\[
\neg\exists x\,\forall y\, L(x,y) \equiv \forall x\,\exists y\,\neg L(x,y)
\]

\medskip
\noindent\textbf{(e)}
\[
\exists x\,\neg L(\text{Lydia},x)
\]

\medskip
\noindent\textbf{(f)}
\[
\exists x\,\forall y\,\neg L(y,x)
\]

\medskip
\noindent\textbf{(g)}

\textit{At least one:}
\[
\exists x\,\forall y\, L(y,x)
\]

\textit{At most one:}
\[
\forall x_1\,\forall x_2\,\Big(\big(\forall y\, L(y,x_1)\wedge \forall y\, L(y,x_2)\big)\Rightarrow x_1=x_2\Big)
\]

\textit{Exactly one} (conjunction of the above):
\[
\exists x\,\forall y\, L(y,x)\ \wedge\ \forall x_1\,\forall x_2\,\Big(\big(\forall y\, L(y,x_1)\wedge \forall y\, L(y,x_2)\big)\Rightarrow x_1=x_2\Big)
\]

Using the uniqueness template $\exists! x\, P(x) \equiv \exists x\big(P(x)\wedge\forall z(P(z)\Rightarrow z=x)\big)$ with $P(x) \equiv \forall y\, L(y,x)$, this simplifies to:
\[
\exists x\,\Big(\forall y\, L(y,x)\ \wedge\ \forall z\big(\forall y\, L(y,z)\Rightarrow z=x\big)\Big)
\]

\medskip
\noindent\textbf{(h)}

\textit{At least two:}
\[
\exists x_1\,\exists x_2\,\big((x_1\neq x_2)\wedge L(\text{Lynn},x_1)\wedge L(\text{Lynn},x_2)\big)
\]

\textit{At most two:}
\[
\forall x_1\,\forall x_2\,\forall x_3\,\Big(\big(L(\text{Lynn},x_1)\wedge L(\text{Lynn},x_2)\wedge L(\text{Lynn},x_3)\big)\Rightarrow(x_1=x_2\vee x_1=x_3\vee x_2=x_3)\Big)
\]

\textit{Exactly two} (combining at least two with at most two, reusing $x_1,x_2$):
\[
\exists x_1\,\exists x_2\,\Big((x_1\neq x_2)\wedge L(\text{Lynn},x_1)\wedge L(\text{Lynn},x_2)
\]
\[
\wedge\ \forall x_3\big(L(\text{Lynn},x_1)\wedge L(\text{Lynn},x_2)\wedge L(\text{Lynn},x_3)\Rightarrow(x_1=x_2\vee x_1=x_3\vee x_2=x_3)\big)\Big)
\]

\medskip
\noindent\textbf{(i)}
\[
\forall x\, L(x,x)
\]

\medskip
\noindent\textbf{(j)}
\[
\exists x\,\forall y\,\big(x\neq y \Rightarrow \neg L(x,y)\big)
\]
\end{solution}

% ======================================================================
\begin{exercise}{13}
    Let $M(x,y)$ be ``$x$ has sent $y$ an e-mail message'' and
    $T(x,y)$ be ``$x$ has telephoned $y$,'' where the domain consists
    of all students in your class. Use quantifiers to express each of
    these statements. (Assume that all e-mail messages that were sent
    are received, which is not the way things often work.)
    \begin{enumerate}[label=(\alph*)]
        \item Chou has never sent an e-mail message to Koko.
        \item Arlene has never sent an e-mail message to or telephoned Sarah.
        \item José has never received an e-mail message from Deborah.
        \item Every student in your class has sent an e-mail message to Ken.
        \item No one in your class has telephoned Nina.
        \item Everyone in your class has either telephoned Avi or sent him an e-mail message.
        \item There is a student in your class who has sent everyone else in your class an e-mail message.
        \item There is someone in your class who has either sent an e-mail message or telephoned everyone else in your class.
        \item There are two different students in your class who have sent each other e-mail messages.
        \item There is a student who has sent himself or herself an e-mail message.
        \item There is a student in your class who has not received an e-mail message from anyone else in the class and who has not been called by any other student in the class.
        \item Every student in the class has either received an e-mail message or received a telephone call from another student in the class.
        \item There are at least two students in your class such that one student has sent the other e-mail and the second student has telephoned the first student.
        \item There are two different students in your class who between them have sent an e-mail message to or telephoned everyone else in the class.
    \end{enumerate}
\end{exercise}

\begin{solution}
Let $M(x,y)$ mean ``$x$ has sent $y$ an e-mail message'' and $T(x,y)$
mean ``$x$ has telephoned $y$.''

\medskip
\noindent\textbf{(a)}
\[
\neg M(\text{Chou},\text{Koko})
\]

\medskip
\noindent\textbf{(b)}
\[
\neg\big(M(\text{Arlene},\text{Sarah})\vee T(\text{Arlene},\text{Sarah})\big)
\]

\medskip
\noindent\textbf{(c)}
\[
\neg M(\text{Deborah},\text{José})
\]

\medskip
\noindent\textbf{(d)}
\[
\forall x\, M(x,\text{Ken})
\]

\medskip
\noindent\textbf{(e)}
\[
\forall x\,\neg T(x,\text{Nina})
\]

\medskip
\noindent\textbf{(f)}
\[
\forall x\,\big(T(x,\text{Avi})\vee M(x,\text{Avi})\big)
\]

\medskip
\noindent\textbf{(g)}
\[
\exists x\,\forall y\,\big(x\neq y \Rightarrow M(x,y)\big)
\]

\medskip
\noindent\textbf{(h)}
\[
\exists x\,\forall y\,\big(x\neq y \Rightarrow (M(x,y)\vee T(x,y))\big)
\]

\medskip
\noindent\textbf{(i)}
\[
\exists x\,\exists y\,\big(x\neq y\wedge M(x,y)\wedge M(y,x)\big)
\]

\medskip
\noindent\textbf{(j)}
\[
\exists x\, M(x,x)
\]

\medskip
\noindent\textbf{(k)}
\[
\exists x\,\forall y\,\big((x\neq y)\Rightarrow(\neg M(y,x)\wedge \neg T(y,x))\big)
\]

\medskip
\noindent\textbf{(l)}
\[
\forall x\,\exists y\,\big(x\neq y\wedge(T(y,x)\vee M(y,x))\big)
\]

\medskip
\noindent\textbf{(m)}
\[
\exists x_1\,\exists x_2\,\big(M(x_1,x_2)\wedge T(x_2,x_1)\big)
\]

\medskip
\noindent\textbf{(n)}
\[
\exists x_1\,\exists x_2\,\forall y\,\Big((x_1\neq x_2\wedge x_1\neq y\wedge x_2\neq y)
\]
\[
\Rightarrow\big(M(x_1,y)\vee M(x_2,y)\vee T(x_1,y)\vee T(x_2,y)\big)\Big)
\]
\end{solution}

% ======================================================================
\begin{exercise}{17}
    Express each of these system specifications using predicates,
    quantifiers, and logical connectives, if necessary.
    \begin{enumerate}[label=(\alph*)]
        \item Every user has access to exactly one mailbox.
        \item There is a process that continues to run during all error conditions only if the kernel is working correctly.
        \item All users on the campus network can access all websites whose url has a .edu extension.
        \item There are exactly two systems that monitor every remote server.
    \end{enumerate}
\end{exercise}

\begin{solution}

\medskip
\noindent\textbf{(a)} Let $A(x,y)$: ``user $x$ has access to mailbox $y$.''
\[
\forall x\,\exists y\,\Big(A(x,y)\wedge\forall z\big(A(x,z)\Rightarrow y=z\big)\Big)
\]

\medskip
\noindent\textbf{(b)} Let $C(x,y)$: ``process $x$ continues to run during error condition $y$'';
$K$: ``the kernel is working correctly.''
\[
\exists x\,\Big(\big(\forall y\, C(x,y)\big)\Rightarrow K\Big)
\]

\medskip
\noindent\textbf{(c)} Let $A(x,y)$: ``user $x$ on the campus can access website $y$'';
$U(y)$: ``website $y$ has a .edu extension.''
\[
\forall x\,\forall y\,\big(U(y)\Rightarrow A(x,y)\big)
\]

\end{solution}

% ======================================================================
\begin{exercise}{21}
    Use predicates, quantifiers, logical connectives, and mathematical
    operators to express the statement that every positive integer is
    the sum of the squares of four integers.
\end{exercise}

\begin{solution}
\[
\forall x\,\Big(x>0\ \Rightarrow\ \exists y_1\,\exists y_2\,\exists y_3\,\exists y_4\,\big(x=y_1^2+y_2^2+y_3^2+y_4^2\big)\Big)
\]
\end{solution}

% ======================================================================
\begin{exercise}{23}
    Express each of these mathematical statements using predicates,
    quantifiers, logical connectives, and mathematical operators.
    \begin{enumerate}[label=(\alph*)]
        \item The product of two negative real numbers is positive.
        \item The difference of a real number and itself is zero.
        \item Every positive real number has exactly two square roots.
        \item A negative real number does not have a square root that is a real number.
    \end{enumerate}
\end{exercise}

% ======================================================================
\begin{exercise}{31}
    Express the negations of each of these statements so that all
    negation symbols immediately precede predicates.
    \begin{enumerate}[label=(\alph*)]
        \item $\forall x\,\exists y\,\forall z\,T(x,y,z)$
        \item $\forall x\,\exists y\,P(x,y)\vee\forall x\,\exists y\,Q(x,y)$
        \item $\forall x\,\exists y\,\big(P(x,y)\wedge\exists z\,R(x,y,z)\big)$
        \item $\forall x\,\exists y\,\big(P(x,y)\to Q(x,y)\big)$
    \end{enumerate}
\end{exercise}

% ======================================================================
\begin{exercise}{37}
    Express each of these statements using quantifiers. Then form the
    negation of the statement so that no negation is to the left of a
    quantifier. Next, express the negation in simple English. (Do not
    simply use the phrase ``It is not the case that.'')
    \begin{enumerate}[label=(\alph*)]
        \item Every student in this class has taken exactly two mathematics classes at this school.
        \item Someone has visited every country in the world except Libya.
        \item No one has climbed every mountain in the Himalayas.
        \item Every movie actor has either been in a movie with Kevin Bacon or has been in a movie with someone who has been in a movie with Kevin Bacon.
    \end{enumerate}
\end{exercise}

% ======================================================================
\begin{exercise}{45}
    Determine the truth value of the statement $\forall x\,\exists y\,(xy=1)$
    if the domain for the variables consists of
    \begin{enumerate}[label=(\alph*)]
        \item the nonzero real numbers.
        \item the nonzero integers.
        \item the positive real numbers.
    \end{enumerate}
\end{exercise}

% ======================================================================
\begin{exercise}{47}
    Show that the two statements $\neg\exists x\,\forall y\,P(x,y)$
    and $\forall x\,\exists y\,\neg P(x,y)$, where both quantifiers
    over the first variable in $P(x,y)$ have the same domain, and
    both quantifiers over the second variable in $P(x,y)$ have the
    same domain, are logically equivalent.
\end{exercise}

% ======================================================================
\begin{exercise}{49}
    \begin{enumerate}[label=(\alph*)]
        \item Show that $\forall x\,P(x)\wedge\exists x\,Q(x)$ is
        logically equivalent to $\forall x\,\exists y\,\big(P(x)\wedge Q(y)\big)$,
        where all quantifiers have the same nonempty domain.
        \item Show that $\forall x\,P(x)\vee\exists x\,Q(x)$ is
        equivalent to $\forall x\,\exists y\,\big(P(x)\vee Q(y)\big)$,
        where all quantifiers have the same nonempty domain.
    \end{enumerate}
\end{exercise}

% ======================================================================
\begin{exercise}{51}
    A statement is in \emph{prenex normal form} (PNF) if and only if
    it is of the form
    \[
    Q_1x_1\,Q_2x_2\cdots Q_kx_k\ P(x_1,x_2,\dots,x_k),
    \]
    where each $Q_i$, $i=1,2,\dots,k$, is either the existential
    quantifier or the universal quantifier, and $P(x_1,\dots,x_k)$ is
    a predicate involving no quantifiers. For example,
    $\exists x\,\forall y\,\big(P(x,y)\wedge Q(y)\big)$ is in prenex
    normal form, whereas $\exists x\,P(x)\vee\forall x\,Q(x)$ is not
    (because the quantifiers do not all occur first).

    Every statement formed from propositional variables, predicates,
    $T$, and $F$ using logical connectives and quantifiers is
    equivalent to a statement in prenex normal form.

    Show how to transform an arbitrary statement to a statement in
    prenex normal form that is equivalent to the given statement.

    \textit{Note: A formal solution of this exercise requires use of
    structural induction, covered in Section 5.3.}
\end{exercise}


\end{document}